\section{Projective Algebraic Sets}

\begin{exercise}
    Let $I$ be a homogeneous ideal in $k[X_1, ..., X_{n+1}]$. Show that $I$ is prime if and only if the following condition is satisfied; for any forms $FG \in k[X_1, ..., X_{n+1}]$, if $FG \in I$, then $F\in I$ or $G \in I$. \\
\end{exercise}

\begin{solution}
    If $I$ is prime, then the condition is clearly satisfied. Suppose now that $I$ satisfies the given condition and let $F,G \in k[X_1, ..., X_{n+1}]$ such that $F = \sum F_i$, $G = \sum G_i$ and $FG \in I$. Since
    $$FG = F_0G_0 + (F_0G_1 + F_1G_0) + (F_0G_2 + F_1G_1 + F_2G_0) + ...,$$
    and $I$ is homogeneous, then each term in the sum above is in $I$. Hence, $F_0G_0 \in I$. Since $I$ satisfies the given condition, then $F_0 \in I$ or $G_0 \in I$. If $F_0 \notin I$, then $G_0 \in I$, and so
    $$F_0G_1 = (F_0G_1 + F_1G_0) - F_1G_0 \in I.$$
    This again implies that $G_1 \in I$ since $F_0 \in I$. By induction, we can show that this case implies that $G_i \in I$ for all $i$, and hence, that $G \in I$. Similarly, if $F_0 \in I$, then
    $$F_1G_0 = (F_0G_1 + F_1G_0) - F_0G_1 \in I$$
    and hence, $F_1 \in I$ or $G_0 \in I$. by repeating this procedure, we can either fall in the case where all the $F_i$'s are in $I$, and hence deduce that $F \in I$, or, there will be one $F_i$ which is not in $I$, and hence, forcing all the $G_i$'s to be in $I$, which would imply that $G \in I$. Thus, in all possible cases, $F \in I$ or $G \in I$. Therefore, $I$ is prime. \\
\end{solution}

\begin{exercise}
    If $I$ is a homogeneous ideal, show that $\rad(I)$ is also homogeneous. \\
\end{exercise}

\begin{solution}
    Let $F \in \rad(I)$ such that $F = \sum_{i=m}^d F_i$ where $\deg(F_i) = i$, then $F^n \in I$ for some $n \geq 1$. It follows that
    $$F^n = F_m^n + [\text{terms of higher degree}] \in I$$
    Since $I$ is homogeneous and $F_m^n$ is the form of smallest degree in the decomposition of $F^n$ into a sum of forms, then $F_m^n \in I$. Hence, $F_m \in \rad(I)$, and so $F - F_m \in \rad(I)$. It suffices to repeat this argument (induction) to get that $F_i \in \rad(I)$ for all $i$. Therefore, $\rad(I)$ is homogeneous. \\
\end{solution}

\begin{exercise}
    State and prove the projective analogues of properties (1)-(10) of Chapter 1, Sections 2 and 3. \\
\end{exercise}

\begin{solution}
    \begin{enumerate}[label = (\arabic*)]
        \item Let's show that if $I$ is the ideal in $k[X_1, ..., X_{n+1}]$ generated by the subset $S \subset k[X_1, ..., X_{n+1}]$, then $V(S) = V(I)$. Let $P \in V(I)$, then $F(P) = 0$ for all $F \in I$. Since $S \subset I$, then $F(P) = 0$ for all $F \in S$. Hence, $V(I) \subset V(S)$. Conversely, let $P \in V(S)$, then $F(P) = 0$ for all $F \in S$. Let $G \in I$, then $G = \sum \alpha_i F_i$ for some $\alpha_i \in k[X_1, ..., X_{n+1}]$ and $F_i \in S$. Let $(x_1, ..., x_{n+1})$ be a choice of homogeneous coordinates for $P$, then
        $$G(x_1, ..., x_{n+1}) = \sum \alpha_i(x_1, ..., x_{n+1}) F_i(x_1, ..., x_{n+1}) = 0.$$
        Since it holds for all choice of homogeneous coordinates for $P$, then $G(P) = 0$. Since it holds for all $G \in I$, then $P \in V(I)$. Therefore, $V(S) = V(I)$.
        \item Let's show that if $\{I_{\alpha}\}$ is any collection of ideals, then $V(\bigcup_{\alpha}I_{\alpha}) = \bigcap_{\alpha}V(I_{\alpha})$. Let $P \in V(\bigcup_{\alpha}I_{\alpha})$, then $F(P) = 0$ for all $F \in \bigcup_{\alpha}I_{\alpha}$. Let $\alpha$ be arbitrary and $F \in I_{\alpha}$, then $F \in \bigcup_{\alpha}I_{\alpha}$, and hence, $F(P) = 0$. Thus, $P \in V(I_{\alpha})$. Since $\alpha$ was set to be arbitrary, then $P \in \bigcap_{\alpha}V(I_{\alpha})$. Conversely, let $P \in \bigcap_{\alpha}V(I_{\alpha})$ and take any $F \in \bigcup_{\alpha}I_{\alpha}$, then $F \in I_{\alpha'}$ for some $\alpha'$. Since $P \in V(I_{\alpha'})$, then $F(P) = 0$. Since it holds for all $F \in \bigcup_{\alpha}I_{\alpha}$, then $P \in V(\bigcup_{\alpha}I_{\alpha})$. Therefore, $V(\bigcup_{\alpha}I_{\alpha}) = \bigcap_{\alpha}V(I_{\alpha})$.
        \item Let's show that if $I \subset J$, then $V(I) \supset V(J)$. Let $P \in V(J)$, then $F(P) = 0$ for all $F \in J$. Let $F \in I$, then $F \in J$, and hence, $F(P) = 0$. Since it holds for all $F \in I$, then $P \in V(I)$. Therefore, $V(I) \supset V(J)$.
        \item Let's show that $V(FG) = V(F)\cup V(G)$, and that $V(I) \cup V(J) = V(\{FG | F\in I, G \in J\})$. Let $P \in V(FG)$, then $(FG)(P) = 0$. Hence, $F(x_1, ..., x_{n+1}) = 0$ or $G(x_1, ..., x_{n+1}) = 0$ for all choice of homogeneous coordinates for $P$. Fix a choice of homogeneous coordinates $P = [x_1 : ... : x_{n+1}]$ and define the polynomials
        $$f(\lambda) = F(\lambda x_1, ..., \lambda x_{n+1}), \qquad g(\lambda) = G(\lambda x_1, ..., \lambda x_{n+1}).$$
        By our previous remark, for all non-zero $\lambda \in k$, either $f(\lambda) = 0$ or $g(\lambda) = 0$. Since $k$ is infinite, then it follows that either $f$ or $g$ has infinitely many zeros, without loss of generality, suppose that $f$ has infinitely many zeros. Since $f$ is a one-variable polynomial, then we must have $f = 0$. Equivalently, this means that $F(P) = 0$. Thus, $P \in V(F) \subset V(F) \cup V(G)$. Conversely, suppose that $P \in V(F)\cup V(G)$ and suppose that $P \in V(F)$ without loss of generality, then $F(P) = 0$. Equivalently, this means that $F(x_1, ..., x_{n+1}) = 0$ for all choice of homogeneous coordinates for $P$. Hence, $F(x_1, ..., x_{n+1})G(x_1, ..., x_{n+1}) = 0$ for all choice of coordinates for $P$. By definition, this implies that $(FG)(P) = 0$, and hence, that $P \in V(FG)$. Therefore, $V(FG) = V(F) \cup V(G)$.\\
        
        Let's now prove the more general statement using what we just proved. Let $P \in V(\{FG | F \in I, G \in J\})$, then $(FG)(P) = 0$ for all $F \in I$ and $G \in J$. By the previous proof, it follows that $F(P) = 0$, or $G(P) = 0$. Now, suppose that $P \in V(I)$, then we are done. Otherwise, there is a $F_0 \in I$ such that $F_0(P) \neq 0$. Thus, for all $G \in J$, the equation $(F_0G)(P) = 0$ implies that $F_0(P) = 0$ or $G(P) = 0$, which implies that $G(P) = 0$. Thus, in that case, $P \in V(J)$. Thus, in both cases, $P \in V(I) \cup V(J)$. Conversely, let $P \in V(I) \cup V(J)$ and suppose without loss of generality that $P \in V(I)$, then $F(P) = 0$ for all $F \in I$. It follows that $(FG)(P) = 0$ for all $F \in I$ and $G \in J$. Therefore, $P \in V(\{FG | F \in I, G \in J\})$. Therefore, $V(I) \cup V(J) = V(\{FG | F\in I, G \in J\})$.
        \item Let's show that $V(0) = \P^n$; $V(1) = \varnothing$; $\{[x_1 : ... : x_{n+1}]\}$ is a projective algebraic set for all $P \in \P^n$. First, we clearly have that every point in $\P^n$ is a zero of the zero polynomial so $V(0) = \P^n$. Similarly, no point in $\P^n$ is a zero of the polynomial 1 so $V(1) = \varnothing$. Next, let $P$ be a point in $\P^n$ and fix some homogeneous coordinates $(x_1, ..., x_{n+1})$ for $P$. In $\A^{n+1}(k)$, the line $L$ through the origin and $(x_1, ..., x_{n+1})$ is an algebraic set, so it can be described as $V(f_1, ..., f_r) \subset \A^{n+1}(k)$ for some polynomials $f_i \in k[x_1, ..., x_{n+1}]$ (Problem 2.15). For all $\lambda \neq 0$, the point $(\lambda x_1, ..., \lambda x_{n+1})$ is in the line, and hence, it is in $V(f_1, ..., f_r)$. Thus, $f_i(\lambda x_1, ..., \lambda x_{n+1}) = 0$ for all $i$ and for all $\lambda \neq 0$. By definition, this means that $P$ is a zero for all $f_i$, so $\{P\} \subset V(f_1, ..., f_r)$ viewed now as a subset of $\P^n$. Conversely, let $Q \in \P^n$ such that $Q \in V(f_1, ..., f_r)$, then choosing some coordinates $(y_1, ..., y_{n+1})$ for $Q$, we get that $(y_1, ..., y_{n+1}) \in V(f_1, ..., f_r) = L \subset \A^{n+1}(k)$. Hence, $(y_1, ..., y_{n+1}) = (\lambda x_1, ..., \lambda x_{n+1})$ for some $\lambda \neq 0$ so $Q = P$. Therefore, $V(f_1, ..., f_r) = \{P\}$ where $V(f_1, ..., f_r)$ is now viewed as a subset of $\P^n$.
        \item Let's prove that if $X \subset Y$, then $I(X) \supset I(Y)$. Let $F \in I(Y)$, then $F(P) = 0$ for all $P \in Y$. Let $P \in X$, then $P \in Y$, and hence, $F(P) = 0$. Since it holds for all $P \in X$, then $F \in I(X)$. Therefore, $I(X) \supset I(Y)$.
        \item Let's show that $I(\varnothing) = k[X_1, ..., X_{n+1}]$ and $I(\P^n(k)) = (0)$. First, for all $F \in k[X_1, ..., X_{n+1}]$, all the points in $\varnothing$ are zeros of $F$, so $F \in I(\varnothing)$. Therefore, $I(\varnothing) = k[X_1, ..., X_{n+1}]$. Next, if $F \in I(\P^n(k))$, then $F(\A^{n+1}(k)\setminus \{0\}) = \{0\}$. If we define $f(\lambda) = F(\lambda, \dots, \lambda)$, we get that $f$ is zero for all $\lambda \neq 0$. Since $f$ is a one-variable polynomial and $k$ is infinite, then $f = 0$, and so $F(0) = 0$. Therefore, $F \in I(\A^{n+1}(k))$, and hence, $F = 0$. It follows that $I(\P^n(k)) = (0)$.
        \item Let's show that $I(V(S)) \supset S$ for any set $S$ of polynomials; $V(I(X)) \supset X$ for any set $X$ of points. Let $F \in S$ and $P \in V(S)$, then by definition of $V(S)$, $F(P) = 0$. Since it holds for all $P \in V(S)$, then $F \in I(V(S))$. Therefore, $I(V(S)) \supset S$. Next, Let $P \in X$ and $F \in I(X)$, then by definition, $F(P) = 0$. Since it holds for all $F \in I(X)$, then $P \in V(I(X))$. Therefore, $V(I(X)) \supset X$.
        \item Let's show that
        $$V(I(V(S))) = V(S) \quad \text{ and } \quad I(V(I(X))) = I(X)$$
        for any set $S$ of polynomials, and any set $X$ of points. First, we know that $V(I(V(S))) \supset V(S)$ and $I(V(S)) \supset S$ from property (8). The inclusion $I(V(S)) \supset S$ implies $V(I(V(S))) \subset V(S)$ using property (3). Therefore, $V(I(V(S))) = V(S)$. Next, we have the inclusions $I(V(I(X))) \supset I(X)$ and $V(I(X)) \supset X$ from property (8). The latter implies that $I(V(I(X))) \subset I(X)$ using property (6). Therefore, $I(V(I(X))) = I(X)$.
        \item Let's show that $I(X)$ is a radical ideal for any $X \subset \P^n(k)$. Let $F^n \in I(X)$, then $F^n(P) = 0$ for all $P \in X$. By property (4), it follows that $F(P) = 0$ or $F^{n-1}(P) = 0$. If $F^{n-1}(P) = 0$, simply reapply the same procedure to get that in all cases, $F(P) = 0$. Since it holds for all $P \in X$, then $F \in I(X)$. Therefore, $I(X)$ is a radical ideal. \\
    \end{enumerate}
\end{solution}

\begin{exercise}
    Show that each irreducible component of a cone is also a cone. \\
\end{exercise}

\begin{solution}
    From the equation $C(V_p(I)) = V_a(I)$, we get that the cone of a projective algebraic subset of $\P^n(k)$ is an affine algebraic set of $\A^{n+1}(k)$. Morever, it is easy to show that $C(V_1 \cup V_2) = C(V_1) \cup C(V_1)$. Hence, if we take a projective algebraic set $V$ such that $V = V_1 \cup \dots \cup V_r$ is its unique decomposition into irreducible components, then $C(V) = C(V_1) \cup \dots \cup C(V_r)$. Let's show that $C(V_1), ..., C(V_r)$ are the irreducible components of $C(V)$. First, they are all irreducible because $I_a(C(V_i)) = I_p(V_i)$ is prime. Morever, suppose that $C(V_i) \subset C(V_j)$, then $I_p(V_i) \subset I_p(V_j)$, and hence, $V_i \subset V_j$. But this is a contradiction since $V_i$ and $V_j$ are irreducible components of $V$. Therefore, $C(V_1), ..., C(V_r)$ are the irreducible components of $C(V)$. By uniqueness of the decomposition into irreducible components, we get that each irreducible component of a cone is also a cone. \\
\end{solution}

\begin{exercise}
    Let $V = \P^1$, $\Gamma_h(V) = k[X,Y]$. Let $t = X/Y \in k_h(V)$, and show that $k(V) = k(t)$. Show that there is a natural one-to-one correspondence between the points of $\P^1$ and the DVR's with quotient field $k(V)$ that contain $k$ (see Problem 2.27); which DVR corresponds to the point at infinity? \\
\end{exercise}

\begin{solution}
    Let $f$ and $g$ be two forms in $\Gamma_h(V)$ of the same degree, say $d$, then we can write $f = \sum_{i=0}^{d}a_i X^{d-i}Y^i$ and $g = \sum_{i=0}^{d}b_iX^{n-i}Y^i$ where $a_i, b_i \in k$, then
    \begin{align*}
        \frac{f}{g} &= \frac{\sum_{i=0}^{d}a_i X^{d-i}Y^i}{\sum_{i=0}^{d}b_i X^{d-i}Y^i} \\
        &= \sum_{i=0}^{d}a_i\frac{X^{d-i}Y^i}{\sum_{j=0}^{d}b_j X^{d-j}Y^j} \\
        &= \sum_{i=0}^{d}a_i\left(\frac{\sum_{j=0}^{d}b_j X^{d-j}Y^j}{X^{d-i}Y^i}\right)^{-1} \\
        &= \sum_{i=0}^{d}a_i\left(\sum_{j=0}^{d}b_j\frac{X^{d-j}Y^j}{X^{d-i}Y^i}\right)^{-1} \\
        &= \sum_{i=0}^{d}a_i\left(\sum_{j=0}^{d}b_j\frac{X^{i-j}}{Y^{i-j}}\right)^{-1} \\
        &= \sum_{i=0}^{d}a_i\left(\sum_{j=0}^{d}b_j t^{i-j}\right)^{-1} \in k(t).\\
    \end{align*}
    Hence, $k(V) \subset k(t)$. Conversely, since $k$ and $t$ are in $k(V)$ and $k(V)$ is a field, then $k(t) \subset k(V)$. Therefore, $k(V) = k(t)$. 

    Next, since $k(V) = k(t)$, then the DVR's with quotient field $k(V)$ that contain $k$ are precisely $\O_a$ for $a \in \A^1(k)$ with uniformizing parameter $t = X - a$, and $\O_{\infty}$ with uniformizing parameter $t = 1/X$. Associate each point $[a : 1] \in \P$ to $\O_a$ and the point at infinity $[0 : 1]$ to $\O_{\infty}$. \\
\end{solution}

\begin{exercise}
    Let $I$ be a homogeneous ideal in $k[X_1, ..., X_{n+1}]$, and
    $$\Gamma = k[X_1, ..., X_{n+1}]/I.$$
    Show that the forms of degree $d$ in $\Gamma$ form a finite-dimensional vector space over $k$. \\
\end{exercise}

\begin{solution}
    \td \\
\end{solution}

\begin{exercise}
    \td \\
\end{exercise}

\begin{solution}
    \td \\
\end{solution}

\begin{exercise}
    \td \\
\end{exercise}

\begin{solution}
    \td \\
\end{solution}

\begin{exercise}
    Let $H_1, ..., H_m$ be hyperplanes in $\P^n$, $m \leq n$. Show that $H_1 \cap H_2 \cap ... \cap H_m \neq \varnothing$. \\
\end{exercise}

\begin{solution}
    Since $H_i$ is a hyperplane, then $H_i = V(f_i)$ where $f_i \in k[X_1, ..., X_{n+1}]$ is a form of degree. Hence, $\bigcap_i H_i = V(f_1, ..., f_m)$. An element $P \in \P^n$ is an element of $\bigcap_i H_i$ if and only if $f_i(P) = 0$ for all $1 \leq i \leq m$. If we write $f_i = a_{i1}X_1 + a_{i2}X_2 + \dots + a_{i(n+1)}X_{n+1}$, then the condition is equivalent to the solvability of the following system of equations:
    \begin{align*}
        a_{11}x_1 + a_{12}x_2 + \dots + a_{1(n+1)}x_{n+1} &= 0, \\
        a_{21}x_1 + a_{22}x_2 + \dots + a_{2(n+1)}x_{n+1} &= 0, \\
        \vdots \qquad \qquad \qquad \qquad \vdots \qquad & \quad \ \vdots \\
        a_{m1}x_1 + a_{m2}x_2 + \dots + a_{m(n+1)}x_{n+1} &= 0.
    \end{align*}
    Define the linear transformation $T : k^{n+1} \to k^m$ as follows:
    $$T(\vec{x}) = \begin{pmatrix}
        a_{11} & a_{12} & \dots & a_{1(n+1)} \\
        a_{21} & a_{22} & \dots & a_{2(n+1)} \\
        \ldots & \ldots & \ddots & \ldots \\
        a_{m1} & a_{m2} & \dots & a_{m(n+1)}
    \end{pmatrix}\vec{x},$$
    then the system of equation is equivalent to the equation $T\vec{x} = 0$ with $x \in k^{n+1}$. By the rank-nullity theorem, we have that
    $$\dim \ker T + \dim \Ima T = n+1.$$
    Since $\Ima T \subset k^m$, then
    $$\dim \ker T = (n+1) - \dim \Ima T \geq (n+1) - m \geq 1.$$
    Thus, there is at least one non-zero vector $(x_1, ..., x_{n+1})$ such that all of its scalar multiples are solution to the equation (by linearity). Hence, if we define $P = [x_1 : ... : x_{n+1}] \in \P^n$, then another way to phrase the previous sentence is to say that $f_i(P) = 0$ for all $i$, or equivalently, that $P \in \bigcap_i H_i$. Therefore, $H_1 \cap H_2 \cap ... \cap H_m \neq \varnothing$. \\
\end{solution}

\begin{exercise}
    \td \\
\end{exercise}

\begin{solution}
    \td \\
\end{solution}

\begin{exercise}
    \td \\
\end{exercise}

\begin{solution}
    \td \\
\end{solution}

\begin{exercise}
    \td \\
\end{exercise}

\begin{solution}
    \td \\
\end{solution}

\begin{exercise}
    \td \\
\end{exercise}

\begin{solution}
    \td \\
\end{solution}

\begin{exercise}
    \td \\
\end{exercise}

\begin{solution}
    \td \\
\end{solution}

\begin{exercise}
    \td \\
\end{exercise}

\begin{solution}
    \td \\
\end{solution}

